% Generated draft: requires mathematical and source review.
\documentclass[11pt]{article}
\usepackage[margin=25mm]{geometry}
\usepackage{fontspec}
\setmainfont{Latin Modern Roman}
\usepackage{amsmath,amssymb,mathtools,unicode-math}
\setmathfont{Latin Modern Math}
\usepackage{enumitem,needspace,xurl,hyperref}
\setlength{\parindent}{0pt}
\setlength{\parskip}{5pt}
\setlength{\emergencystretch}{3em}
\newcommand{\N}{\mathbb{N}}
\newcommand{\Z}{\mathbb{Z}}
\newcommand{\Q}{\mathbb{Q}}
\newcommand{\R}{\mathbb{R}}
\newcommand{\C}{\mathbb{C}}
\newcommand{\F}{\mathbb{F}}
\newcommand{\A}{\mathbb{A}}
\newcommand{\PP}{\mathbb{P}}
\newcommand{\E}{\mathbb{E}}
\newcommand{\eps}{\varepsilon}
\newcommand{\dd}{\,\mathrm{d}}
\begin{document}
\section*{UnsolvedMath ID: 30006614}
\textbf{Source problem code:} catalogue entry
\subsection{Definitions and mathematical statement}
For a set \(X\), a split is an unordered pair \(\{H,X\setminus H\}\) of nonempty complementary subsets. Half-space membership data record whether an observed point belongs to a specified side of a split. The supplied record does not specify the underlying set, the family of splits, the observed incidences, or the meaning and extent of noise.

For a metric space \((M,d)\), define the metric interval by
\[
I_d(a,b)=\{u\in M:d(a,u)+d(u,b)=d(a,b)\}.
\]
The space is a median metric space if, for every \(a,b,c\in M\),
\[
I_d(a,b)\cap I_d(b,c)\cap I_d(c,a)
\]
consists of exactly one point, called the median of \(a,b,c\). A connected graph is a median graph if its vertex set, equipped with the shortest-path metric in which every edge has length one, is a median metric space. In a graph, a half-space conventionally means a nonempty proper vertex subset such that both it and its complement are geodesically convex; the record does not specify an ambient graph for its data.

\textbf{Source problem, with unresolved parameters.} Construct a coarse analogue of the Buneman graph for missing or noisy half-space membership data, representing the available partial data and permitting computation of medians with controlled error. Determine hypotheses under which the resulting metric space has quantitatively controlled median geometry.

This is a construction and characterization problem schema, rather than a fully specified mathematical assertion. The source does not determine a Buneman-graph convention, an admissible input class, a representation criterion, an error model, an approximate-median condition, or the parameters on which the quantitative bounds may depend. Consequently, no particular universal quantifier over inputs or numerical error inequality is asserted here. In particular, the exact median definition above supplies terminology and does not assert that the requested coarse model is an exact median space.

The extracted original statement also contains an incomplete question about being countably infinite and a question about whether a minimum, after choosing a peak, might lie on the boundary and be well-defined. Their objects and relationship to this problem are unspecified. These fragments cannot be converted into cardinality, boundary, or existence hypotheses. The subsequent text explicitly introduces Problem~6 on graphs of hyperplanes and does not provide a reliable continuation of this problem.
\subsection{Short English statement}
The supplied problem asks for a Buneman-type model of incomplete or noisy half-space observations, together with median computations having quantitative error bounds and conditions ensuring controlled median geometry. The input model and the meaning of these bounds are unspecified, so this remains a qualified problem schema requiring source clarification and human mathematical review.
\subsection{Sources}
\begin{itemize}
\item Local UnsolvedMath registry, supplied complete record: permanent ID \(30006614\), title \emph{Coarse Buneman Graphs for Partial Data}, exact source code \texttt{catalogue entry}, in the Oberwolfach Reports Open Problems collection. The supplied proposed year is 2026; no individual proposer is identified.
\item Supplied report citation: \emph{Median Geometry and Applications} (2026), Oberwolfach Reports, pp.~487--540, DOI \texttt{10.4171/owr/2026/8}. Supplied source link: \url{https://doi.org/10.4171/owr/2026/8}. The background describes these reports as documenting mathematical workshops and their open-problem discussions.
\item Supplied literature source, labelled ``OWR report'': \url{https://ems.press/content/serial-article-files/53603}.
\item The supplied literature assessment, checked on \textbf{2026-08-22}, identifies the intended question as Problem~5, describes it as open, and reports that no resolving result was found. This is supplied dated curation, not an independently verified literature assessment; time-sensitive status and cited claims require independent verification.
\item The registry labels the statement \texttt{corrected\_verified} and states that only clear formula delimiters, notation, grammar, and typography were repaired, preserving mathematical content and scope. This supplied editorial qualification does not resolve the missing formulation details or the contamination in the extracted original statement.
\end{itemize}
\end{document}
